% !TeX root = ../navier-stokes-manuscript.tex
% ============================================================
% 2.1 Vortex Stretching
% ============================================================

The momentum equation has four pieces:
\[
  \underbrace{\partial_t u}_{\text{local acceleration}}
  +
  \underbrace{(u\cdot\nabla)u}_{\text{nonlinear transport}}
  =
  \underbrace{-\nabla p}_{\text{pressure}}
  +
  \underbrace{\nu\Delta u}_{\text{viscous diffusion}}.
\]
The regularity pressure sits between nonlinear transport and viscous diffusion.
The nonlinear term uses the velocity to carry and deform its own field, while viscosity spreads spatial variation through the fluid.
The hope is simple: viscosity smooths the fluid faster than the nonlinear motion can sharpen it.
If that remains true at every scale reached by the flow, the velocity stays smooth.

\subsection{Vortex Stretching}\label{sub:ThePerfectlyStretchingVortex}

The local rotation of the velocity is \(\omega:=\nabla\times u\), and the part of the velocity gradient that stretches or contracts material directions is
\[
  S:=\frac12\bigl(\nabla u+(\nabla u)^{\mathsf T}\bigr),
\]
Taking curl of the momentum equation removes pressure.
For an incompressible velocity field, nonlinear transport leaves the vorticity-stretching action \(S\omega\).
When the vorticity points along a positive strain direction, \(S\omega\) points with it and amplifies the rotation already present in the fluid.
A rotating region organized along such a direction is the material vortex segment we can now follow.

Take one such segment, with axial direction \(e\) and length \(L(t)\).
On an interval where the vorticity and the positive strain direction remain aligned, write
\[
  Se=\sigma(t)e,
  \qquad
  \sigma(t)>0,
  \qquad
  \frac{\dot L(t)}{L(t)}
  =
  e\cdot Se
  =
  \sigma(t).
\]
The segment lengthens along \(e\).
Incompressibility fixes what happens across it at the same time.
Let \(\mathcal V\) be its preserved material volume, set \(A(t):=\mathcal V/L(t)\), and write \(r_{\mathrm{eff}}(t):=\sqrt{A(t)/\pi}\).
Then
\[
  \nabla\cdot u=0,
  \qquad
  \frac{d}{dt}(A L)=0,
  \qquad
  \frac{\dot A}{A}=-\sigma(t),
  \qquad
  \frac{\dot r_{\mathrm{eff}}}{r_{\mathrm{eff}}}
  =
  -\frac{\sigma(t)}{2}.
\]
The same material segment becomes longer and thinner.

The axial strain also acts on the rotation carried by that narrowing segment.
The viscous vorticity equation is
\[
  \frac{D\omega}{Dt}
  =
  S\omega+\nu\Delta\omega,
  \qquad
  \frac{D}{Dt}
  :=
  \partial_t+u\cdot\nabla,
\]
and when \(\omega=|\omega|e\),
\[
  S\omega
  =
  \sigma(t)\omega,
  \qquad
  \frac12\frac{D|\omega|^2}{Dt}
  =
  \sigma(t)|\omega|^2
  +
  \nu\,\omega\cdot\Delta\omega.
\]
Viscosity remains present in this identity.
On any interval where the right-hand side is positive, the rotation strengthens while the radius falls.
Vorticity is the rotating part of the velocity gradient, so this growth is already gradient growth.
If \(|u|\leq U_0\) on the material segment \(\Omega_{\mathrm{seg}}(t)\), then
\[
  E_{\mathrm{seg}}(t)
  :=
  \frac12\int_{\Omega_{\mathrm{seg}}(t)}|u(x,t)|^2\,dx
  \leq
  \frac12U_0^2\mathcal V
  <
  \infty,
  \qquad
  \left|\frac12\bigl(\nabla u-(\nabla u)^{\mathsf T}\bigr)\right|_{\mathrm F}
  =
  \frac{|\omega|}{\sqrt2}
  \leq
  |\nabla u|_{\mathrm F}.
\]
The segment can carry finite kinetic energy while its radius falls and the rotating part of its velocity gradient grows.

\divider{\NavierStokes{} Scaling}

That coexistence changes the original viscosity hope.
The sharpening is now concentrated across a smaller width, so the four momentum terms must be compared at a finer scale to see whether viscosity gains relative strength there.
Fix \(t_0\) and set \(\ell:=r_{\mathrm{eff}}(t_0)\).
For \(\lambda>1\), define
\[
  u_\lambda(x,t)
  :=
  \lambda u(\lambda x,\lambda^2t),
  \qquad
  p_\lambda(x,t)
  :=
  \lambda^2p(\lambda x,\lambda^2t).
\]
This is a scale comparison, not a later state of the material segment.
At time \(\lambda^{-2}t_0\), the vortex in \(u_\lambda\) has width \(\lambda^{-1}\ell\), and the four terms become
\[
\begin{aligned}
  \partial_tu_\lambda(x,t)
  &=
  \lambda^3(\partial_tu)(\lambda x,\lambda^2t),
  \\
  (u_\lambda\cdot\nabla)u_\lambda(x,t)
  &=
  \lambda^3\bigl((u\cdot\nabla)u\bigr)(\lambda x,\lambda^2t),
  \\
  \nabla p_\lambda(x,t)
  &=
  \lambda^3(\nabla p)(\lambda x,\lambda^2t),
  \\
  \nu\Delta u_\lambda(x,t)
  &=
  \lambda^3\nu(\Delta u)(\lambda x,\lambda^2t).
\end{aligned}
\]
Viscosity strengthens at the finer scale, but nonlinear transport strengthens by the same \(\lambda^3\) factor; local acceleration and pressure remain in that balance as well.
Going smaller has not given viscosity an automatic advantage.
The balance survives the change of scale, leaving the size used to measure concentration as the next distinction.

\divider{Critical Height}

The homogeneous Sobolev scale places kinetic energy and velocity gradients on one continuum.
Let \(\Lambda:=(-\Delta)^{1/2}\).
A change of variables gives
\[
  \norm{u_\lambda(t)}_{\dot H^s}^2
  =
  \lambda^{2s-1}
  \norm{u(\lambda^2t)}_{\dot H^s}^2.
\]
The exponent vanishes at \(s=\tfrac12\).
Set
\[
  H(t)
  :=
  \frac12\norm{u(t)}_{\dot H^{1/2}}^2
  =
  \frac12\norm{\Lambda^{1/2}u(t)}_2^2.
\]
With \(H_\lambda\) denoting the critical height of \(u_\lambda\),
\[
  \norm{u_\lambda(t)}_2^2
  =
  \lambda^{-1}\norm{u(\lambda^2t)}_2^2,
  \qquad
  H_\lambda(t)=H(\lambda^2t),
  \qquad
  \norm{\nabla u_\lambda(t)}_2^2
  =
  \lambda\norm{\nabla u(\lambda^2t)}_2^2.
\]
Across the finer comparison, kinetic energy falls while the first-derivative energy rises.
The critical height remains fixed.
Finite energy can miss the concentration because it lies below this invariant height.

\divider{Nonlinear Production versus Viscous Dissipation}

The invariant height now has to be read along the original solution.
Apply \(\Lambda^{1/2}\) to the momentum equation and pair with \(\Lambda^{1/2}u\).
The signed nonlinear production is
\[
  \mathcal P_{1/2}(t)
  :=
  -\operatorname{Re}
  \pair
    {\Lambda^{1/2}u(t)}
    {\Lambda^{1/2}\bigl((u(t)\cdot\nabla)u(t)\bigr)}_{L^2}.
\]
The pressure pairing vanishes because \(\nabla\cdot u=0\), while the Laplacian contributes \(-\nu\norm{\Lambda^{3/2}u}_2^2\).
The critical height obeys
\[
  H'(t)
  +
  \nu\norm{\Lambda^{3/2}u(t)}_2^2
  =
  \mathcal P_{1/2}(t).
\]
It rises exactly when nonlinear production exceeds viscous dissipation.
Under the same scale comparison, the two contributions satisfy
\[
  \mathcal P_{1/2}[u_\lambda](t)
  =
  \lambda^2\mathcal P_{1/2}[u](\lambda^2t),
  \qquad
  \nu\norm{\Lambda^{3/2}u_\lambda(t)}_2^2
  =
  \lambda^2\nu\norm{\Lambda^{3/2}u(\lambda^2t)}_2^2.
\]
Production and dissipation remain matched in scaling.
Their common \(\lambda^2\) factor is the inverse of the time contraction \(t\mapsto\lambda^{-2}t\).

\divider{The Heat Clock}

That inverse time factor is the clock viscosity assigns to a width.
For the linear heat equation \(v_t=\nu\Delta v\),
\[
  \widehat v(\xi,t+\delta t)
  =
  e^{-\nu|\xi|^2\delta t}\widehat v(\xi,t).
\]
A scale-localized structure of width \(r\) occupies frequencies \(|\xi|\simeq r^{-1}\).
Viscosity changes that scale by order one when \(\nu|\xi|^2\delta t\simeq1\), which gives the comparison time
\[
  \tau_\nu(r)
  :=
  \frac{r^2}{\nu}.
\]
At the finer width,
\[
  \tau_\nu(\lambda^{-1}r)
  =
  \lambda^{-2}\tau_\nu(r),
\]
The same change of scale that narrows the structure also shortens its viscous comparison clock by the square of the spatial factor.

\divider{The Zeno Cascade}

One contraction does not create a cascade.
It gives a width--time relation whose repetition can be tested as a candidate schedule.
For the geometric sequence
\[
  r_n:=2^{-n}r_0,
  \qquad
  \tau_n:=\frac{r_n^2}{\nu},
\]
the heat clocks satisfy
\[
  \tau_n
  =
  \frac{r_0^2}{\nu}4^{-n},
  \qquad
  \sum_{n=0}^{\infty}\tau_n
  =
  \frac{4r_0^2}{3\nu}
  <
  \infty.
\]
The finite sum proves only that these comparison clocks do not rule out accumulation at a finite time.
It does not produce a singular solution.
For the sequence to describe a candidate singularity, one \NavierStokes{} solution would have to pass through widths
\[
  r_0>r_1>r_2>\cdots\longrightarrow0
\]
at times
\[
  t_0<t_1<t_2<\cdots\uparrow T_*<\infty,
  \qquad
  t_{n+1}-t_n\asymp\tau_n,
\]
with comparison constants independent of \(n\).
The remaining time would then satisfy
\[
  T_*-t_n
  \asymp
  \sum_{k=n}^{\infty}\tau_k
  =
  \frac{4r_n^2}{3\nu}.
\]
Each finer width would have to form inside an interval tending to zero at order \(\tau_n\).
Under this candidate schedule, the velocity would have to change on each shorter interval.
Its local acceleration carries the first rate of that change, its local jerk carries the next, and the demand continues through every later temporal response.
The spatial sequence has become the demand that opens the temporal Tower.
